% Chapter 2: Kinds of Sentences

\chapter{Kinds of Sentences}

\section{Pedagogical Purpose}
This chapter builds on the understanding of complete sentences by introducing how sentences can be used for different purposes. It helps learners communicate more effectively and understand the intentions behind different types of spoken and written language. It also reinforces the importance of correct punctuation.

\begin{learningobjectives}
The learner can:
\begin{itemize}
    \item identify statements, questions, commands, and exclamations.
    \item explain the purpose of each kind of sentence.
    \item use appropriate punctuation (full stop, question mark, exclamation mark) for each sentence kind.
    \item transform one kind of sentence into another (e.g., statement to question).
    \item construct sentences of different kinds to express specific intentions.
\end{itemize}
\end{learningobjectives}

\section{Prerequisite Check}
\begin{quickcheck}
Read the following groups of words. Put a tick (✓) next to the ones that are complete sentences and add a full stop.
\begin{enumerate}
    \item The children played in the park [✓].
    \item Running very fast [ ].
    \item My friend Maya is reading a book [✓].
    \item Ate breakfast quickly [ ].
    \item The sun shines brightly [✓].
\end{enumerate}
\end{quickcheck}

\section{Chapter Opener}
\subsection*{A Busy Morning}
\textbf{Ms. Sharma:} Good morning, Class! I have some exciting news today.

\textbf{Rohan:} Really? What is it?

\textbf{Maya:} Please tell us, Ms. Sharma!

\textbf{Ms. Sharma:} We are going on a field trip next month!

\textbf{Leo (the parrot, squawking loudly):} "A field trip! Wow!"

\textbf{Ms. Sharma:} Yes, Leo, it's wonderful news! Now, everyone, please open your grammar books to Chapter 2.

\section{Language Experience}
Read this short conversation between Maya and Rohan.

\textbf{Maya:} Did you finish your homework?
\textbf{Rohan:} Yes, I did. I wrote about my favourite animal.
\textbf{Maya:} That's great! My favourite animal is a dolphin.
\textbf{Rohan:} Tell me more about dolphins.
\textbf{Maya:} They are very intelligent creatures!
\textbf{Rohan:} Wow, that's amazing!

\section{Look and Notice}
\begin{enumerate}
    \item Which sentences ask a question? What punctuation mark do they end with?
    \item Which sentences tell you something? What punctuation mark do they end with?
    \item Which sentences show strong feeling or excitement? What punctuation mark do they end with?
    \item Can you find a sentence that gives an instruction or makes a request?
\end{enumerate}

\section{Discover / Explicit Teaching}
\subsection*{Guided Discovery: Sentences for Different Jobs}
In our daily lives, we use sentences for many different reasons. Sometimes we want to tell someone something, sometimes we want to ask, sometimes we want to give an order, and sometimes we want to show how we feel!

Let's look at the sentences from our opener and language experience:
\begin{itemize}
    \item "We are going on a field trip next month!" (Tells us something, shows excitement)
    \item "What is it?" (Asks something)
    \item "Please tell us, Ms. Sharma!" (Makes a request)
    \item "They are very intelligent creatures!" (Tells us something, shows strong feeling)
    \item "Please open your grammar books to Chapter 2." (Gives an instruction)
\end{itemize}
Notice how the punctuation at the end changes, and how the words are arranged.

\section{Core Explanation}
Every time we speak or write, we use sentences to do different jobs. There are four main kinds of sentences:

\begin{enumerate}
    \item \textbf{Statements (Declarative Sentences):}
    \begin{itemize}
        \item \textbf{Purpose:} To tell you something or give information.
        \item \textbf{Punctuation:} Ends with a full stop (.).
        \item \textbf{Example:} Maya loves to read books.
    \end{itemize}

    \item \textbf{Questions (Interrogative Sentences):}
    \begin{itemize}
        \item \textbf{Purpose:} To ask something.
        \item \textbf{Punctuation:} Ends with a question mark (?).
        \item \textbf{Example:} Does Rohan play football?
    \end{itemize}

    \item \textbf{Commands (Imperative Sentences):}
    \begin{itemize}
        \item \textbf{Purpose:} To give an order, make a request, or give an instruction. The subject (you) is often understood but not stated.
        \item \textbf{Punctuation:} Ends with a full stop (.) or sometimes an exclamation mark (!) if the command is strong.
        \item \textbf{Example:} Please close the door. (Request)
        \item \textbf{Example:} Close the door! (Strong command)
    \end{itemize}

    \item \textbf{Exclamations (Exclamatory Sentences):}
    \begin{itemize}
        \item \textbf{Purpose:} To show strong feelings like surprise, excitement, anger, or joy.
        \item \textbf{Punctuation:} Ends with an exclamation mark (!).
        \item \textbf{Example:} What a beautiful day it is!
    \end{itemize}
\end{enumerate}

\section{Guided Practice}
\begin{activitybox}{Activity A: Identify the Kind of Sentence}
Write S for Statement, Q for Question, C for Command, or E for Exclamation.
\begin{enumerate}
    \item The birds are singing outside. [S]
    \item What is your favourite colour? [Q]
    \item Please sit down quietly. [C]
    \item How beautiful the flowers are! [E]
    \item Have you seen Maya's new book? [Q]
\end{enumerate}
\end{activitybox}

\section{Independent Practice}
\begin{activitybox}{Activity B: Construct Sentences}
Write one sentence of each kind about your school.
\begin{enumerate}
    \item Statement: \underline{\hspace{10cm}}
    \item Question: \underline{\hspace{10cm}}
    \item Command: \underline{\hspace{10cm}}
    \item Exclamation: \underline{\hspace{10cm}}
\end{enumerate}
\end{activitybox}

\section{Summary}
\begin{summarybox}
    \begin{itemize}
        \item A \textbf{statement} tells something and ends with a full stop (.).
        \item A \textbf{question} asks something and ends with a question mark (?).
        \item A \textbf{command} gives an order or request and ends with a full stop (.) or an exclamation mark (!).
        \item An \textbf{exclamation} shows strong feeling and ends with an exclamation mark (!).
    \end{itemize}
\end{summarybox}